\section{投资问题：一体化究竟保护了什么}

一体化常被简化为“原料自给、成本更低”，但对恒逸石化而言，更准确的描述是：文莱一期提供成品油与芳烃利润池，国内PTA和聚酯承接原料转化，CPL--PA6形成另一条尼龙链，供应链服务支持周转；不同环节在不同周期承担利润或库存。核心判断是，一体化可以降低单一环节断供和利润池完全外流的风险，却不能消除油价、价差、库存、少数股东与资本开支波动。

2025年涤纶收入597.44亿元，是集团最大收入板块；炼油和化工产品合计收入395.68亿元，是文莱利润弹性的主要载体；PTA、PIA和锦纶规模相对较小。收入体量与盈利弹性并不相同：聚酯贡献最大收入，却只有4.07\%毛利率；文莱炼油和芳烃收入较小，但2026H1的利润跃升主要来自该端价差与满负荷。

\begin{exhibitbox}[2025年分部收入、毛利与研究含义]
\noindent\begin{tabularx}{\textwidth}{L{2.55cm}R{2.0cm}R{1.7cm}R{2.0cm}X}
\toprule
\textbf{板块} & \textbf{收入} & \textbf{毛利率} & \textbf{毛利约} & \textbf{研究含义} \\
\midrule
炼油产品 & 284.78亿元 & 4.86\% & 13.84亿元 & 2025低基数；2026价差弹性最大 \\
化工产品 & 110.90亿元 & 7.63\% & 8.46亿元 & PX／苯为主要芳烃价值池 \\
PTA & 44.16亿元 & 0.37\% & 0.16亿元 & 产能过剩下接近盈亏边缘 \\
PIA & 14.12亿元 & 4.55\% & 0.64亿元 & 小体量差异化品种 \\
涤纶 & 597.44亿元 & 4.07\% & 24.34亿元 & 按年报收入减成本；体量最大，受加工费和库存约束 \\
锦纶及副产品 & 18.20亿元 & 6.83\% & 1.24亿元 & 新项目爬坡，尚非核心利润 \\
供应链服务 & 65.67亿元 & 3.91\% & 2.57亿元 & 低毛利周转业务，不给独立估值溢价 \\
\bottomrule
\end{tabularx}
\end{exhibitbox}
\sourcenote{HY-OFF-2025AR，2025年度分产品收入与毛利率；毛利为收入乘毛利率的近似，财务截止2025-12-31。}

这张分部表解释了2026年反转为何不是简单收入增长：利润主要由文莱炼油／芳烃毛利率跃迁和国内低毛利板块边际修复驱动。2026E收入1,417.4亿元只比2025A高约24.8\%，而归母利润从2.58亿元升至77.5亿元，真正的变量是价差、负荷与利润归属。

\section{文莱与国内平台是两套经济模型}

文莱一期面对东南亚成品油和芳烃市场，原油采购以Brent为基准、主要来自文莱和马来西亚；产品销往东南亚和澳大利亚。其利润由综合裂解价差、负荷、原油贴水、运费、美元融资、税惠和70\%权益共同决定。国内平台面对中国PTA--聚酯过剩、织造开工与出口需求，利润更接近加工费和库存周期。把两者用同一“油价上涨利好／利空”逻辑解释，会遗漏利润池在链条中的迁移。

\begin{exhibitbox}[两套经济模型的差异]
\small
\noindent\begin{tabularx}{\textwidth}{L{2.35cm}L{3.7cm}L{3.7cm}X}
\toprule
\textbf{维度} & \textbf{文莱炼化／芳烃} & \textbf{国内PTA／聚酯／尼龙} & \textbf{估值含义} \\
\midrule
定价 & 原油、成品油、PX／苯美元定价 & PX、PTA、MEG、长丝及人民币／出口价格 & 汇率和区域价差需分开 \\
需求 & 东南亚／澳大利亚油品与芳烃 & 中国织造、服装、家纺、包装与出口 & 需求锚不可互相替代 \\
盈利代理 & 综合裂解、PX--石脑油、负荷 & PTA／聚酯加工费、库存、开工 & 不能用单一油价方向推利润 \\
归属 & 70\%控股，30\%少数股东 & 控股与参股口径混合 & 总利润不等于归母 \\
现金约束 & 原油库存、美元融资、运费与税 & 库存销售、应收／预付、内外销与信用期 & OCF是共同验证项 \\
成长路径 & 一期优化、二期规划 & 技改、瓶片／再生、CPL--PA6爬坡 & 未披露项目不给基准信用 \\
\bottomrule
\end{tabularx}
\end{exhibitbox}
\sourcenote{HY-OFF-2025AR；HY-OFF-P2；IND-BRUNEI；IND-CZCE；研究截止2026-07-22。}

投资含义是，文莱高景气能够暂时覆盖国内链条低毛利，却也让集团盈利更依赖单一海外利润池。若文莱价差回落而国内加工费未接棒，一体化只会把利润重新分散到低回报资产，不会自动保持2026H1水平。

\section{产能地图：设计能力不是可归母销量}

公司披露文莱一期原油加工能力800万吨／年、成品油565万吨／年、PX与苯合计265万吨／年；国内参控股PTA产能2,150万吨／年、PIA 30万吨／年、PET Fiber 938万吨／年、Bottle Chip含RPET 530万吨／年，以及CPL／PA6 100／60万吨／年。该地图展示业务广度，但不同数字包含控股、参股和产品口径，不能直接相加成“总产能”。

\begin{exhibitbox}[资产与产能地图]
\small
\noindent\begin{tabularx}{\textwidth}{L{2.55cm}R{2.1cm}L{2.6cm}L{2.3cm}X}
\toprule
\textbf{资产块} & \textbf{设计／参控股产能} & \textbf{2025产销锚} & \textbf{归属／口径} & \textbf{关键缺口} \\
\midrule
文莱一期原油加工 & 800万吨／年 & 炼油产603.64、销606.95万吨 & 恒逸文莱70\%控股 & 实际进料、综合收率与分产品利润 \\
成品油 & 565万吨／年 & 包含在炼油产品口径 & 文莱产品 & 分产品销量、ASP与客户 \\
PX／苯 & 合计265万吨／年 & 化工产219.27、销217.96万吨 & 文莱产品 & PX／苯拆分价差与利润 \\
PTA & 2,150万吨／年 & 控股口径产220.91、销222.94万吨 & 参控股混合 & 权益、开工、内部消化和外销量 \\
PIA & 30万吨／年 & 产20.81、销21.38万吨 & 直接板块 & 客户、ASP与成本拆分 \\
PET Fiber & 938万吨／年 & 涤纶产867.97、销868.93万吨 & 产品范围与产能口径未完全一致 & 分品种量价、开工和库存 \\
Bottle Chip／RPET & 530万吨／年 & 未单列 & 参控股口径 & 收入、认证、客户与毛利 \\
CPL／PA6 & 100／60万吨／年 & 锦纶产8.89、销4.67万吨 & 新增制造板块 & 爬坡、良率、认证与盈亏平衡 \\
\bottomrule
\end{tabularx}
\end{exhibitbox}
\sourcenote{HY-OFF-2025AR；analysis/company\_fundamental\_cards.md，数据截止2025-12-31，研究截止2026-07-22。PTA与聚酯产能含参控股口径，禁止直接相加或全部并表。}

产能地图的“所以呢”是：本报告只用已披露产销量和分部收入检查模型，不使用设计产能乘行业加工费。特别是PTA 2,150万吨／年与控股口径222.94万吨销量差异巨大，若不处理权益和内部消化，利润可能被高估数倍。

\section{供应链服务与内部协同：周转价值大于估值溢价}

供应链服务收入65.67亿元、毛利率3.91\%，其经济作用更接近采购、销售、物流和周转支撑，而不是独立高回报业务。关联交易中还存在原油采购、PTA采购、聚酯采购、货运和工程管理。它们提高原料与产品组织效率，却可能扩大营运资金需求；因此供应链服务应与库存、预付款、应收款和票据一起评价。

若集团通过内部平台集中采购获得更优贴水、缩短物流或降低库存，价值最终应表现为文莱或国内分部毛利与现金周转改善，而不是供应链收入规模本身。反之，若交易额增长伴随存货和预付款上升，所谓协同可能只增加资产负债表占用。

\begin{keyinsight}[商业模式的“所以呢”]
恒逸的护城河不是“拥有最多产能”，而是能够在海外炼化与国内聚酯之间组织原料、产品和资本。\textbf{行动上，只有当这一组织能力转化为可归母利润和经营现金流，才值得提高估值；产能、收入和关联交易规模本身不增加目标价。}
\end{keyinsight}

\section{一体化的反方情景}

最强反方情景是国内链条长期过剩导致PTA与聚酯资产回报低迷，而文莱高景气又因区域供给恢复快速回落；此时完整链条反而需要更多库存、担保和资本开支，利润池迁移不能覆盖资本成本。2025年PTA毛利率0.37\%、涤纶4.07\%、文莱主体净利润2.42亿元，已经展示这类低回报组合并非理论风险。

证明反方情景错误，需要看到三项证据同时改善：文莱在正常化价差下仍产生可观现金利润；国内PTA和聚酯在库存下降后提升加工费而非只增销量；资本开支和营运资金不再吞噬经营现金流。当前H1预告只覆盖第一项的高景气方向，尚未完成全部证明。

\begin{riskbox}[商业模式失效阈值]
若文莱分部利润快速回落、国内PTA加工费重返约250元／吨以下、长丝库存继续上升，同时H1经营现金流仍大额为负，则“一体化缓冲”无法支撑7.5倍基准PE。\textbf{行动上，应同时下调H2利润和估值倍数，而不是用另一环节的设计产能填补缺口。}
\end{riskbox}
